\documentclass[11pt,a4paper]{article}
\usepackage[T1]{fontenc}
\usepackage[utf8]{inputenc}
\usepackage{lmodern}
\usepackage{amsmath,amssymb,amsthm,mathtools,mathrsfs}
\usepackage{graphicx,float,xcolor,multicol}
\usepackage[shortlabels]{enumitem}
\usepackage[margin=27mm]{geometry}
\usepackage{microtype,needspace,placeins}
\usepackage{tikz,tikz-cd}
\usetikzlibrary{arrows.meta,calc,positioning,patterns,decorations.pathreplacing}
\usepackage[colorlinks=true,linkcolor=teal,urlcolor=teal,citecolor=teal]{hyperref}
\setlength{\parindent}{0pt}
\setlength{\parskip}{.6em}
\setcounter{secnumdepth}{0}
\allowdisplaybreaks
\raggedbottom
\providecommand{\cross}{\times}
\DeclareUnicodeCharacter{2020}{\ensuremath{\dagger}}
\DeclareMathOperator{\supp}{supp}
\DeclareMathOperator*{\esssup}{ess\,sup}
\providecommand{\chapter}{\section}
\newenvironment{tool}[2]{\subsection*{Tool #1 --- #2}}{}
\newcommand{\ToolRef}[1]{Tool #1}
\newcommand{\FigureTag}[2]{\textit{Figure #1}}
\newcommand{\Result}[1]{\subsection*{#1}}
\newcommand{\Application}[1]{\subsubsection*{Application\if\relax\detokenize{#1}\relax\else\space--- #1\fi}}
\newcommand{\ProofHeading}{\par\textit{Proof.}\space}
\newcommand{\ProofOf}[1]{\par\textit{Proof of #1.}\space}
\newcommand{\RemarkHeading}{\par\textbf{Remark.}\space}
\newcommand{\NoteHeading}{\par\textbf{Note.}\space}
\newcommand{\ExerciseHeading}{\par\textbf{Exercise.}\space}

\graphicspath{{figures/complex-analysis/}}
\begin{document}
\section*{The Complex Plane}
% BLOG-CONTENT-BEGIN
% Overleaf publishing connection verified.

\subsection*{Tool 1 — Basis-Fixing Reduction}
Consider a linear transformation,
\[
T\left(\sum_{i=1}^{n}a_i e_i\right)=\sum_{i=1}^{n}a_iT(e_i).
\]
Observe that if $T$ fixes the basis $\{e_i\}$, then $T$ is the identity map. Thus, compose $T$ with linear transformations so that the composition fixes $\{e_i\}$; then one can trace the steps back to recover $T$.

A simple application of the aforementioned tool goes into showing that any isometry $T\colon\mathbb C\to\mathbb C$ is of the form
\[
T(z)=z_0+\omega z
\qquad\text{or}\qquad
T(z)=z_0+\omega\overline z.
\]


\textit{Proof.}
Without loss of generality, let $T$ fix $0$; otherwise consider $T(z)-T(0)$. Then
\begin{align*}
|T(v)-T(0)|
&=\operatorname{dist}(T(v),T(0))
 =\operatorname{dist}(v,0)\\
&=|v-0|=|v|,
\end{align*}
so $|T(v)|=|v|$. Now,
\begin{align*}
2\langle x,y\rangle
&=|x|^2+|y|^2-|x-y|^2\\
&=|T(x)|^2+|T(y)|^2-
  \operatorname{dist}(T(x),T(y))^2\\
&=|T(x)|^2+|T(y)|^2-|T(x)-T(y)|^2\\
&=2\langle T(x),T(y)\rangle.
\end{align*}
Hence $\{T(e_i)\}_i$ is orthogonal, and therefore forms a basis. If
\[
T(v)=\sum_i b_iT(e_i)=T\left(\sum_i a_ie_i\right),
\]
then, noting that
\[
\langle T(v),T(e_i)\rangle=\langle v,e_i\rangle,
\]
we get $a_i=b_i$ for every $i$, $1\leq i\leq n$. We know that
\[
T\left(\sum_{i=1}^{n}a_ie_i\right)=\sum_{i=1}^{n}a_iT(e_i).
\]
We apply Tool $1$ by composing $T$ with a rotation $\varphi$ so that
\[
(\varphi\circ T)(e_1)=e_1.
\]
The map $\varphi\circ T$ is again an isometry that preserves $0$ (observe here that $T$ is linear). Thus
\[
\langle(\varphi\circ T)(e_2),e_1\rangle
=\langle(\varphi\circ T)(e_2),(\varphi\circ T)(e_1)\rangle
 =\langle e_2,e_1\rangle=0, \text{ and,}
\]
\[
\langle(\varphi\circ T)(e_2),(\varphi\circ T)(e_2)\rangle
=\langle e_2,e_2\rangle=1.
\]
Therefore \((\varphi\circ T)(e_2)=\pm e_2,\)
and hence
\[
(\varphi\circ T)(z)=z\qquad\text{or}\qquad
(\varphi\circ T)(z)=\overline z.
\]
Consequently,
\[
T(z)=\omega z\qquad\text{or}\qquad T(z)=\omega\overline z,
\qquad \omega\in S^1.
\]
Let $z_0=T(0)$. We obtain the desired forms
\[
T(z)=z_0+\omega z
\qquad\text{or}\qquad
T(z)=z_0+\omega\overline z.
\]

\textbf{Remark.}
The power of Tool $1$ goes far beyond this. Maps that fix certain elements often turn out to be ``nice''. We shall see another application of this soon. Before that, it is worth emphasizing the rigidity of the complex plane: if two points of the plane are connected by a rubber band, there are essentially two ways to deform it preserving lengths -- reflections and rotations. This rigidity feature can be extended to all Euclidean spaces $\mathbb R^n$. However, there are instances of non-Euclidean spaces, for which the doubling map $z\mapsto 2z$ is an isometry. 

\subsubsection*{The dihedral group}
Now let us illustrate another application of Tool $1$ by proving that the dihedral group \(D_n\) is isomorphic to \(\operatorname{Aut}(C_n),\) where $C_n$ is the cyclic graph on $n$ vertices.

\textit{Proof.}
It is easy to check that $D_n$ acts on $C_n$, since rotations and reflections preserve edges. Thus there is a group homomorphism
\[
\varphi\colon D_n\longrightarrow\operatorname{Aut}(C_n).
\]
The map $\varphi$ is clearly injective. We have to prove that it is also surjective. To that end, label the vertices $\{1,2,3,4,\ldots,n\}$. Without loss of generality, for a given $T\in\operatorname{Aut}(C_n)$ assume $T(1)=1$; otherwise compose it with $\varphi(T')$, where $T'$ is a rotation. Write the cycle as
\[
1\longrightarrow2\longrightarrow3\longrightarrow\cdots
\longrightarrow n\longrightarrow1.
\]
If $T(1)=1$, then $T(2)$ must be adjacent to $1$, so $T(2)=2$ or $n$. If $T(2)=2$, then $T(3)$ must be adjacent to $2$ and cannot be $1$. Inductively, $T(i)=i$ for every $i$. If $T(2)=n$, then $T(3)$ must be adjacent to $n$ but not to $1$. Again, inductively, we conclude that $T$ is a reflection. Now, tracing back, we conclude that $T\in\varphi(D_n)$. Hence $\varphi$ is an isomorphism.
% BLOG-CONTENT-END
\end{document}
